\clearpage
\phantomsection% 
\addcontentsline{toc}{chapter}{ABSTRACT}
\thispagestyle{fancy}

\begin{center}
	\large \bfseries \MakeUppercase{Abstract}\\
%	\normalsize \normalfont {\thetitleEN}\\
%	\normalsize \normalfont {\theauthor}\\
	\bigskip
	
	\normalsize \normalfont \justifying \singlespacing
	Student admission is a crucial process in higher education management that requires a structured and efficient system. HKBP Nommensen University currently relies on manual processes in carrying out student admissions, which causes various obstacles such as slow processing of applicant data, difficulties in document verification, and lack of information transparency for prospective students. This research aims to develop an integrated web-based student admission system capable of managing the entire admission process digitally, from registration, data verification, payment, to re-enrollment. The development method used is the Prototype method, chosen because it allows iterative development with direct feedback from users. The system is designed using the Model-View-Controller (MVC) architecture with the Spring Boot framework, MySQL database, and a responsive interface using HTML, CSS, and Bootstrap. The system design includes functional and non-functional requirements analysis, creation of use case diagrams, activity diagrams, sequence diagrams, class diagrams, Data Flow Diagrams (DFD), Entity Relationship Diagrams (ERD), and interface design using Figma. System testing is conducted through three approaches: Black-box Testing to validate functionality, Integration Testing to examine inter-module integration, and System Usability Scale (SUS) to measure the system usability level from the user perspective. The research is expected to produce an efficient, integrated, and user-friendly student admission system for all stakeholders at HKBP Nommensen University.\par
    
    \vspace{20pt}
	\raggedright\textbf{Keywords: student admission, integrated system, prototype, Spring Boot, web}
	
	\vfill
	
\end{center}
\clearpage
